%% related_work.tex
%% Related Work section for the in-progress writeup
%%   "A Systems Characterization of a Multi-Agent Meeting Assistant Pipeline"
%% This file is self-contained and can be compiled on its own or included in
%% the full paper draft. Bibliography entries are collected in references.bib
%% (also in this directory).

\section{Related Work}
\label{sec:related}

\subsection{LLM Inference Serving}

\paragraph{Continuous batching and PagedAttention.}
Orca~\cite{yu2022orca} introduced iteration-level scheduling for transformer
inference, replacing static batching with dynamic request interleaving to
improve GPU utilization.  vLLM~\cite{kwon2023vllm} built on this with
\emph{PagedAttention}, which manages the KV-cache as a set of fixed-size
non-contiguous blocks analogous to OS virtual memory pages.  This allows
multiple requests to share physical KV-cache storage and enables the
prefix-caching feature~\cite{vllm2024prefixcache} that our work
instruments and measures.  Our benchmark characterizes how vLLM's
prefix-cache reuse behaves under multi-agent concurrent dispatch: when
three workers share an identical 66K-token transcript prefix, parallel
dispatch raises server-side cache hit rates by 13–26 percentage points
relative to sequential dispatch.

\paragraph{Prefill/decode disaggregation and chunked prefill.}
Splitwise~\cite{patel2024splitwise} and DistServe~\cite{zhong2024distserve}
separate prefill and decode into distinct serving phases, observing that
their memory-bandwidth and compute profiles differ substantially.
Sarathi-Serve~\cite{agrawal2024sarathiserve} addresses the \emph{prefill
stall} problem by splitting long prefill operations into fixed-size chunks
interleaved with ongoing decode work.  We propose applying
chunked prefill to our workload: the three 66K-token concurrent prefills in
parallel mode drive peak KV-cache occupancy to 59\%, a scaling risk that
chunked prefill directly mitigates by smoothing occupancy over time.

\paragraph{Prefix-caching and shared-context serving.}
SGLang~\cite{zheng2024sglang} introduces \emph{RadixAttention}, which
maintains a radix tree of cached KV blocks indexed by token sequence prefix.
This allows multiple requests sharing a long common prefix (e.g., a system
prompt or document) to reuse previously computed KV blocks without
re-encoding.  Our shared-prefix design in \texttt{worker\_util.py} is
motivated by the same principle: all three worker agents receive the
identical \texttt{(brief + transcript)} prefix so that vLLM's block-level
prefix cache can amortize transcript encoding across concurrent workers.
Parrot~\cite{lin2024parrot} characterizes LLM applications at the semantic
variable level, observing that many production applications submit requests
with large shared prefixes and proposes a semantic caching layer above the
serving engine.  Our measurements validate their observation in the
multi-agent setting: server-side prefix-cache hit rates reach 43–72\%
across our six transcripts.

\subsection{Multi-Agent LLM Systems}

\paragraph{Agent frameworks.}
LangChain~\cite{langchain2023} and its stateful extension
LangGraph~\cite{langgraph2024} provide graph-based orchestration of LLM
agents with first-class support for conditional edges and cyclic flows,
enabling the reviewer reflection loop we implement.  AutoGen~\cite{wu2023autogen}
introduces a conversational multi-agent framework where agents exchange
messages to collaboratively complete tasks, also supporting human-in-the-loop
and code-execution patterns.  MetaGPT~\cite{hong2023metagpt} assigns
software-engineering roles (product manager, architect, engineer) to distinct
agents that coordinate through structured deliverables, a design similar to
our coordinator-worker decomposition.  Our contribution relative to these
frameworks is empirical: we instrument the inference engine beneath the
framework layer and characterize how the orchestration topology (sequential
vs.\ parallel vs.\ LangGraph fan-out) affects KV-cache pressure, TTFT,
and throughput -- measurements invisible to the frameworks themselves.

\paragraph{Reflection and self-correction.}
Reflexion~\cite{shinn2023reflexion} enables agents to reason verbally about
past failures and refine their outputs over multiple trials.
Self-Refine~\cite{madaan2023selfrefine} iteratively improves LLM outputs
using the same model as both generator and critic.  Our bounded reflection
loop (MAX\_REVISIONS = 2) extends this pattern to the multi-agent setting,
where the Reviewer agent cross-checks all three worker outputs against the
original transcript and selectively re-dispatches only failing workers.
Our measurements reveal a key systems consequence of this pattern:
the reflection loop is the dominant source of end-to-end latency variance,
capable of doubling agent-call counts on short transcripts.

\subsection{LLM Application Characterization}

AlpaServe~\cite{li2023alpaserve} profiles model-parallel LLM inference
across multiple GPUs, exposing how pipeline and tensor parallelism interact
with batching efficiency.  FlexGen~\cite{sheng2023flexgen} characterizes
throughput-latency trade-offs when offloading KV-cache to CPU/disk for
high-throughput generation.  Closest to our work, Parrot~\cite{lin2024parrot}
specifically studies \emph{applications} (rather than single requests),
showing that application-level patterns -- shared prefixes, request
dependencies, result reuse -- are invisible to serving engines that reason
only at the request level.  Our paper contributes a concrete data point to
this line of work: a five-agent pipeline with a shared long-context prefix
and a cyclic revision dependency that creates non-trivial inference
engine pressure.  We show that characterizing this workload requires
server-level telemetry (KV-cache occupancy, queue depth, per-request TTFT)
rather than application-layer latency alone.

\subsection{Meeting Summarization}

AMI~\cite{mccowan2005ami} and ICSI~\cite{janin2003icsi} are established
meeting corpora used as benchmarks for abstractive meeting summarization.
QMSum~\cite{zhong2021qmsum} extends this with query-based summarization over
meetings from diverse domains such as parliamentary committee sessions, and
MeetingBank~\cite{hu2023meetingbank} provides long-form city council meeting
transcripts.  These corpora motivate the domain flavour and the size tiers of
our benchmark; to keep the artifact self-contained and free of any third-party
or personal data, the transcripts we ship are \emph{synthetic} stand-ins,
procedurally generated to match those length tiers rather than text drawn from
the corpora themselves.  Hierarchical
summarization approaches~\cite{zhang2022hierarchical} decompose long
documents into segments before summarizing, analogous to our coordinator-
worker decomposition where the Coordinator emits a compact brief that
focuses downstream workers on salient portions of the transcript.

%%---------------------------------------------------------------------
%% Bibliography entries live in references.bib (same directory).
%% Compile with, e.g.:  \bibliographystyle{plain} \bibliography{references}
%%---------------------------------------------------------------------
